\documentclass[11pt]{article}

\usepackage[T1]{fontenc}
\usepackage[utf8]{inputenc}
\usepackage{amsmath,amssymb,mathtools,mathrsfs}
\usepackage[margin=1in]{geometry}
\usepackage{microtype}

\setlength{\parindent}{0pt}
\setlength{\parskip}{0.72em}

\begin{document}

\section*{How the Proof Works}

For the global three-dimensional Navier--Stokes regularity problem, the fluid obeys
\[
\partial_tu+(u\cdot\nabla)u+\nabla p=\nu\Delta u,
\qquad \nabla\cdot u=0,
\qquad u(\cdot,0)=u_0.
\]
A strain in the flow pulls a vortex tube along its axis. As the vortex stretches, incompressibility forces it to contract across its width. Each material parcel keeps its volume, so lengthening the tube narrows its cross-section and the rotating core inside it. That strain also acts on the vorticity \(\omega=\nabla\times u\). Curling the equation gives
\[
\partial_t\omega+(u\cdot\nabla)\omega
=(\omega\cdot\nabla)u+\nu\Delta\omega.
\]
The left side follows the change in rotation with the moving fluid. When the strain is aligned with the vortex, \(\omega\cdot((\omega\cdot\nabla)u)>0\), so the stretching contribution raises \(|\omega|\) while the core narrows. Since \(\omega\) is one spatial derivative of \(u\), its amplification means that velocity changes more sharply across the core. The gradients rise.

As that core narrows, nonlinear transport can sharpen its velocity profile more intensely, but viscosity also acts more strongly across the shorter distance. Scaling says exactly how both effects change when the spatial pattern is compressed. For \(\lambda>1\), set
\[
u_\lambda(x,t)=\lambda u(\lambda x,\lambda^2t),
\qquad
p_\lambda(x,t)=\lambda^2p(\lambda x,\lambda^2t).
\]
This sends a feature of width \(r\), velocity scale \(U\), and duration \(\tau\) to one of width \(r/\lambda\), velocity scale \(\lambda U\), and duration \(\tau/\lambda^2\). Substitution gives
\[
\partial_tu_\lambda+(u_\lambda\cdot\nabla)u_\lambda
+\nabla p_\lambda-\nu\Delta u_\lambda
=\lambda^3
\bigl[\partial_tu+(u\cdot\nabla)u+\nabla p-\nu\Delta u\bigr]
(\lambda x,\lambda^2t)=0.
\]
Every term acquires \(\lambda^3\). The core can get narrower and faster, yet nonlinear transport and viscosity strengthen in exact proportion. Compression does not break their balance.

Energy sees something different. The core loses volume faster than its velocity grows:
\[
\|u_\lambda(\cdot,t)\|_2^2
=\lambda^{-1}\|u(\cdot,\lambda^2t)\|_2^2.
\]
So finite kinetic energy can survive even while the velocity profile steepens. On the homogeneous Sobolev scale,
\[
\|u_\lambda(\cdot,t)\|_{\dot H^s}
=\lambda^{s-\frac12}
\|u(\cdot,\lambda^2t)\|_{\dot H^s}.
\]
Below \(s=\tfrac12\), concentration at smaller scales makes the norm decrease. Above \(s=\tfrac12\), it makes the norm increase. At \(s=\tfrac12\), the factor is one. This is the critical height: narrowing the core does not automatically change its measured size. Write
\[
E_s(t):=\frac12\|\Lambda^su(t)\|_2^2,
\qquad
H(t):=E_{1/2}(t),
\]
for that critical measurement.

The time scaling makes the danger concrete. If each stage contracts by a fixed factor \(0<\rho<1\), then its scales read
\[
r_j=\rho^jr_0,
\qquad
U_j=\rho^{-j}U_0,
\qquad
\tau_j=\rho^{2j}\tau_0.
\]
Their total duration is
\[
\sum_{j=0}^{\infty}\tau_j
=\tau_0\sum_{j=0}^{\infty}\rho^{2j}
=\frac{\tau_0}{1-\rho^2}<\infty.
\]
The widths collapse, the velocity scale diverges, and infinitely many narrower, faster stages fit inside a finite interval. Energy can remain finite while smaller and smaller regions carry steeper gradients over shorter and shorter intervals. This is the Zeno cascade. Scaling does not prove that a solution must follow this pattern. It shows why finite energy and the scale balance between nonlinearity and viscosity do not rule it out.

Now suppose, for contradiction, that the maximal smooth solution ends at a finite time \(T_*\).

In order for a singularity to form, velocity has to blow up. In order for that to happen, its acceleration would also have to blow up, which means its jerk has to blow up, and so on, and so on:
\[
u,\qquad \partial_tu,\qquad \partial_t^2u,\qquad \ldots
\]
What you would expect to see in this pattern, if a singularity were to form, is less spatial regularity as you go up the time derivatives, not more. Navier--Stokes reverses that expectation.

The first three differentiated equations make that reversal visible:
\[
\partial_tu
=\nu\Delta u-(u\cdot\nabla)u-\nabla p,
\]
\[
\partial_t^2u
=\nu\Delta\partial_tu
-(\partial_tu\cdot\nabla)u
-(u\cdot\nabla)\partial_tu
-\nabla\partial_tp,
\]
\[
\begin{aligned}
\partial_t^3u={}&\nu\Delta\partial_t^2u
-(\partial_t^2u\cdot\nabla)u
-2(\partial_tu\cdot\nabla)\partial_tu \\
&-(u\cdot\nabla)\partial_t^2u
-\nabla\partial_t^2p.
\end{aligned}
\]
The complete pattern is the binomial form
\[
\begin{aligned}
\partial_t^{\,n+1}u={}&\nu\Delta\partial_t^n u \\
&-\sum_{a=0}^{n}\binom{n}{a}
\bigl(\partial_t^a u\cdot\nabla\bigr)
\partial_t^{\,n-a}u
-\nabla\partial_t^n p.
\end{aligned}
\]
Only now set
\[
V_n:=\partial_t^n u,
\qquad
Q_n:=\partial_t^n p.
\]
Viscosity has one exact role here: \(\nu\Delta V_n\) connects the next temporal response to two higher spatial derivatives inside the preceding response. The transport terms expand with the product rule and pressure remains attached to every row, but the Laplacian supplies the upward spatial step.

At critical height, the energy calculation translates that connection into Sobolev order:
\[
H,\qquad H',\qquad H'',\qquad H''',\qquad\ldots
\]
expose
\[
E_{1/2},\qquad E_{3/2},\qquad E_{5/2},\qquad E_{7/2},\qquad\ldots
\]
The transport and pressure terms travel with that relation, but they are not the point of the passage.
As the temporal order rises, the Sobolev height rises with it:
\[
\text{higher temporal order}
\quad\longrightarrow\quad
\text{higher spatial Sobolev order}.
\]
If every finite height in that rising sequence is controlled, its payoff is immediate. In three dimensions the Sobolev embedding theorem gives
\[
s>k+\frac32
\quad\Longrightarrow\quad
H^s(\mathbb R^3)\hookrightarrow C^k(\mathbb R^3),
\]
so
\[
\bigcap_{s<\infty}H^s
=H^\infty
\hookrightarrow C^\infty.
\]
The temporal sequence that looked as though it should expose worsening spatial roughness points in the opposite direction: every finite temporal order opens a higher spatial Sobolev order, and all finite spatial orders together give smoothness.

The identities expose that structure; the estimates still have to control it. Each temporal response carries a spatial Sobolev ladder, and one coordinate means one concrete statement:
\[
V_r=\partial_t^ru
\qquad\text{measured in}\qquad
H_x^m.
\]
A mark in this grid identifies the coordinate under discussion. It does not announce a bound. Any finite analytical demand uses only finitely many temporal orders and finitely many spatial heights, so fix those two cutoffs and take the resulting rectangle. On the putative interval \((0,T_*)\), the datum-generated whole-terminal rectangle theorem supplies that control. For every finite \(N,m\),
\[
V_N\in L^2(0,T_*;H^{m+1}),
\qquad
V_{N+1}\in L^2(0,T_*;H^{m-1}).
\]
The constants may depend on \(N\) and \(m\); the proof never asks for one constant covering the infinite grid. The adjacent rows give the temporal trace relation, the whole-terminal estimate controls the finite block, and pressure stays coupled to its velocity coordinates.

These rectangles overlap. Their shared coordinates are derivatives of the velocity history generated by \(u_0\), so the overlaps agree. Compatibility can then assemble all finite blocks into
\[
u_0
\longrightarrow
\{\mathcal R_{N,M}[u_0]\}_{N,M<\infty}
\longrightarrow
\mathcal I[u_0]
\subset
\bigcap_{N,M<\infty}H_t^N(0,T_*;H_x^M).
\]
Projecting this intersection onto its zeroth temporal row gives
\[
\mathscr S_0(0,T_*)
:=
\bigcap_{M<\infty}L^2(0,T_*;H_x^M).
\]
That is the full movement:
\[
\begin{gathered}
\text{identities expose higher heights}
\longrightarrow
\text{rectangles control finite blocks} \\
\longrightarrow
\text{compatibility forms the intersection}.
\end{gathered}
\]

This is where the reversal becomes an endpoint statement. For every finite \(m\), the intersection contains the adjacent pair
\[
u\in L^2(0,T_*;H^{m+1}),
\qquad
u_t\in L^2(0,T_*;H^{m-1}).
\]
The Hilbert-space trace theorem gives
\[
u\in C([0,T_*];H^m)
\qquad\text{for every finite }m.
\]
Because these traces all come from one velocity field, they meet at
\[
u_*:=u(T_*)\in\bigcap_{m<\infty}H^m=H^\infty.
\]
The pressure equation and the differentiated Navier--Stokes recurrence then place every temporal response and its pressure response in \(H^\infty\) as well. Smooth local existence from \(u_*\), followed by uniqueness, continues the velocity history past \(T_*\). A finite maximal time cannot remain, so \(T_*=\infty\).

The rectangle bounds produce \(\mathcal I[u_0]\). What remains is to take the familiar estimates from that intersection.

This is not an a priori estimate proof. It is an intersection-to-estimate proof. Viscosity relates the velocity's time derivatives to its spatial Sobolev derivatives. Taken through every finite order, those two directions meet in a common intersection. At that intersection, \(H^\infty(\mathbb R^3)\hookrightarrow C^\infty(\mathbb R^3)\) gives smoothness, and the familiar estimates can be taken from whichever finite Sobolev height they require. The estimates come from the intersection, not the other way around.

For any finite \(r,k\) and \(2\leq p\leq\infty\), choose
\[
m>k+\frac32-\frac3p.
\]
Sobolev embedding then gives the requested finite norm of \(\nabla^k\partial_t^ru\). The usual examples fall out at once:
\[
\mathcal I[u_0]
\longrightarrow C_tH_x^1
\longrightarrow L_t^\infty L_x^6,
\]
\[
\mathcal I[u_0]
\longrightarrow C_tH_x^2
\longrightarrow L_t^\infty L_x^\infty,
\]
\[
\mathcal I[u_0]
\longrightarrow C_tH_x^3
\longrightarrow L_t^\infty W_x^{1,\infty}.
\]
No single estimate is the destination. Any finite downstream demand selects the spatial height it needs from the intersection. The proof's arc can now be read in one line:
\[
\begin{aligned}
u_0
&\longrightarrow \{(V_n,Q_n)\}_{n\geq0}
\longrightarrow \{\mathcal R_{N,M}[u_0]\}_{N,M<\infty} \\
&\longrightarrow \mathcal I[u_0]
\longrightarrow \mathscr S_0
\longrightarrow u_*\in H^\infty
\longrightarrow T_*=\infty.
\end{aligned}
\]

\end{document}
